\textit{Metamensagem do designer.} Este é o meu entendimento, como designer, de quem você é: uma pessoa que amplia a tela para ler, lê com fluência quando a palavra cabe na vista e pode perder a relação entre um prato e o seu preço. Você quer pedir a sua refeição e conferir a de quem vai comer com você, de preferência comparando antes de decidir e corrigindo antes de registrar, porque um engano só aparece quando a comida chega. Este, portanto, é o sistema que projetei para você: cada escolha reúne nome, opções e preço; você pode voltar a eles, conferir o que já incluiu e corrigir o pedido antes de confirmá-lo; e, quando uma opção não estiver clara, há uma explicação no ponto em que a dúvida surgiu. A interface funciona como preposto do designer \cite[p.~84]{desouza}, e a perspectiva comunicativa admite manipulação direta \cite[p.~230]{barbosa}. A Tabela~\ref{tab:jakobson} decide a equivalência de canal entre os dois receptores.
\begin{table*}[!t]
\caption{Espaço de design com dois receptores (elementos de Jakobson na Engenharia Semiótica)}
\label{tab:jakobson}
\centering\scriptsize\setlength{\tabcolsep}{3pt}\renewcommand{\arraystretch}{1.1}
\begin{tabularx}{\textwidth}{@{}>{\raggedright\arraybackslash}p{.075\textwidth}>{\raggedright\arraybackslash}p{.17\textwidth}>{\raggedright\arraybackslash}X>{\raggedright\arraybackslash}X>{\raggedright\arraybackslash}p{.2\textwidth}@{}}
\hline \textbf{Elemento} & \textbf{Pergunta do designer} & \textbf{Para a persona} & \textbf{Para o agente} & \textbf{Diferença declarada e status}\\ \hline
Emissor & Que limitações e intenções do designer devem ser comunicadas? & O designer se manifesta em preços, opções, respostas e na ajuda; o restaurante é fictício e não processa pagamento. & O mesmo preposto produz a observação seguinte, com os mesmos textos e respostas. & Nenhuma diferença de autoria. PD: autoria comum.\\
Receptor & Que aspectos do usuário, como interpretado pelo designer, devem ser comunicados? & Lívia amplia a tela, pode perder a relação entre prato e preço e quer conferir antes de registrar. & O modelo escolhe uma ação a partir do recorte recebido e da memória limitada de trechos vistos. & O agente não tem a experiência de Lívia. HO: representação funcional da persona.\\
Contexto & Que elementos do contexto o sistema processa, e como? & Reconhece a ocasião (almoço, pedido da irmã, jantar) e lê o combinado na instrução. & Recebe a mesma instrução literal da tarefa. & Ruído, pressa e companhia ficam fora da sessão. HO.\\
Código & Qual é a linguagem de interface? & Português, moeda em reais, etiquetas V/SL/SA, posição, tipografia, cor e controles. & Texto dos fragmentos legíveis, tipo e coordenadas dos elementos. & O agente perde parte das pistas de tipografia e cor. HO.\\
Canal & Por quais canais a metacomunicação chega? & Tela do smartphone vista por uma região ampliada que ela desloca. & Interseções legíveis entre signos e recorte $360/z\times800/z$, com memória FIFO de fragmentos. & Aproximação geométrica da vista ampliada. HO.\\
Mensagem & O que se quer contar, e com que efeito? & Escolher, conferir, corrigir e registrar uma refeição; a confirmação comunica o fim. & A mesma mensagem, recebida pelos mesmos signos e pela mesma confirmação. & O sucesso é decidido pelo verificador externo, nunca pela fala do agente. PD.\\ \hline
\end{tabularx}
\end{table*}

